\section{Introduction}
\label{sec:intro}

Large language models trained with Reinforcement Learning with Verifiable Rewards (RLVR) achieve strong performance on mathematical and logical reasoning~\citep{deepseek2025r1,shao2024grpo}, but their parameter counts and extended chain-of-thought traces make deployment costly.
Structured pruning---removing entire attention heads or MLP channels---offers a hardware-friendly compression path that reduces both memory and latency without specialized sparse-computation support~\citep{llmpruner2023,graspprune2026}.
However, structured pruning at moderate sparsity severely degrades reasoning ability, necessitating a post-pruning recovery stage.
\citet{selfhealing2025} show that RLVR-based recovery substantially outperforms SFT-based recovery for pruned reasoning models, establishing RLVR as the preferred recovery paradigm.

Despite this, current approaches treat pruning and recovery as independent steps: pruning methods optimize for immediate post-pruning accuracy~\citep{resp2025,graspprune2026,bonsai2025}, while recovery methods operate on whatever architecture they receive~\citep{selfhealing2025,paser2025}.
No prior work has studied whether the \emph{choice} of what to prune affects how well RLVR can subsequently recover reasoning---or whether this interaction differs from SFT-based recovery.
This gap is consequential because RLVR and SFT impose fundamentally different functional demands on the retained architecture: \citet{rlsqueezes2026} show that RLVR recovers through global exploration and path compression, while SFT recovers through local trajectory imitation.
The optimal set of structures to retain should therefore depend on the recovery paradigm.

We hypothesize that \emph{functional overlap}---the degree to which a structure's computation is redundant with its within-layer peers---positively predicts RLVR recoverability.
If a pruned structure has high overlap with retained structures, RLVR can absorb its function through global policy optimization over the remaining architecture.
Conversely, structures with unique, non-redundant computations are difficult to recover once removed.
This yields what we term the \textbf{Redundancy Paradox}: traditional pruning preferentially removes high-overlap structures because their function is ``already covered,'' while recovery-aware pruning \emph{retains} them because their function is ``easily redistributed'' by RLVR.
The same structural property receives opposite treatment depending on whether one optimizes for immediate accuracy or post-recovery performance.

We propose \textbf{Recovery-Aware Structured Pruning (\rasp{})}, a three-stage framework that incorporates recoverability into the pruning decision.
In Stage~1, \rasp{} computes for each structure both a gradient-based importance score and a functional overlap score via intra-layer Centered Kernel Alignment (CKA;~\citealp{kornblith2019cka}).
In Stage~2, these are combined into a recovery-aware score $S(h) = I(h) \times (1 - F(h))^\alpha$ that protects structures which are both important \emph{and} functionally unique, while allowing removal of important-but-redundant structures whose function RLVR can redistribute.
In Stage~3, the pruned model undergoes a brief SFT warmup to restore generation competence, followed by GRPO-based RLVR recovery~\citep{shao2024grpo} to restore reasoning ability.
Unlike prior CKA-based pruning that uses inter-layer similarity for depth pruning~\citep{mpruner2024,lachance2024cka}, \rasp{} uses \emph{intra-layer} CKA between individual GQA groups and MLP channels to predict per-structure recoverability---a finer granularity optimizing for a fundamentally different objective.

We evaluate \rasp{} on DeepSeek-R1-Distill-Qwen-7B~\citep{deepseek2025r1} at 20\% structured sparsity.
At matched sparsity, \rasp{}-pruned models retain 29\% GSM8K~\citep{cobbe2021gsm8k} accuracy immediately after pruning, compared to 1\% for importance-only pruning ($N{=}100$, $p{<}0.0001$)---a 28-percentage-point advantage attributable to the recovery-aware scoring criterion.
After SFT warmup, both pruning strategies recover substantially (57--64\%), with detailed comparison deferred to \S\ref{sec:experiments}.
[TBD: GRPO results]

Our contributions are:
\begin{enumerate}[nosep,leftmargin=*]
    \item We identify the \emph{Redundancy Paradox}---that functional redundancy, typically a signal for removal in standard pruning, is the property that enables RLVR recovery---and formalize it through intra-layer CKA-based functional overlap scoring.
    \item We propose \rasp{}, a recovery-aware pruning framework that integrates functional overlap into structural pruning decisions, achieving 29\% post-prune GSM8K accuracy versus 1\% for importance-only pruning at matched 20\% sparsity on DeepSeek-R1-Distill-Qwen-7B.
    \item [TBD: GRPO recovery results---contribution demonstrating that \rasp{}'s post-prune advantage translates to superior post-RLVR reasoning performance.]
    \item We present preliminary evidence that the optimal pruning strategy may be conditioned on the recovery paradigm~\citep{frankle2019lth}, supporting the notion of \emph{paradigm-conditioned prunability} [TBD: pending GRPO confirmation and statistical validation].
\end{enumerate}
